% Einheit 1 von 2 – Der Regelsatz (bank/integrationsregeln/e1.jsonl, zusammenbau v0.1)
%% TODO Zeitmarke Einheit 1: Marken-Zeile nicht lesbar
%% TODO Zweigzeile Teil 1: was hier gelernt wird (Satz in Schülersprache aus dem Einheitstitel) – Einheit 1
\einheitenkopf[e1]{Einheit 1 von 2 · Der Regelsatz}
\zweigzeile{keine Prüfungsaufgabe}
\setcounter{aufgabe}{5}

%% TODO Titel: Ich-kann-Satz fehlt in der Bank (Platzhalter: Kettenname) – Nr. 6
%% TODO Anweisung: ein Satz über der ersten Teilaufgabe fehlt in der Bank – Nr. 6
\begin{aufgabe}{Der Regelsatz}
%% TODO \rechenplatz{n}: Schrittzahl des Grundfalls fehlt in der Bank (nur bei fester Schreibform, 2.3 b)
\begin{teile}
\teil Welche Regel braucht man für $x^7$?\\ \kreuz{Potenzregel}\\ \kreuz{Faktorregel}\\ \kreuz{Konstantenregel}\\ \kreuz{lineare innere Funktion}
\teil Gib eine Stammfunktion an: $f(x) = x^4$ $F(x) = $ \leerfeld
\teil Gib eine Stammfunktion an: $f(x) = 6x^2 - 4x + 3$ $F(x) = $ \leerfeld
\teil Gib eine Stammfunktion an: $f(x) = e^{4x}$ $F(x) = $ \leerfeld
\teil Gib eine Stammfunktion an: $f(x) = 3x^2 + 4x - 2 + e^{2x}$ $F(x) = $ \leerfeld
\end{teile}
\end{aufgabe}
